\section{Taxonomy - IMPROVED VERSION}
\subsection{Definitions}

Building on the formalization from \textit{[1]}, we define \textit{ambiguous} and \emph{unanswerable} natural language questions in the context of text-to-SQL.

Let \textit{D} be a database and \textit{S} the set of all \textit{non-equivalent valid} SQL queries that can be formulated on \textit{D}\footnote{Two queries are considered non-equivalent if there exists some database state on which they produce different results, following the criterion in [1]. Queries that are syntactically different but semantically equivalent (i.e., produce identical results on all possible database states) are treated as a single element in \textit{S}.} and \textit{q} be a natural language question.

Let \textit{K} be an external knowledge base, \(\textit{K} = \{k_1, k_2, \ldots, k_n\}\), where each \(k_i\) is a piece of evidence that might assist in interpreting \textit{q} and mapping it to a valid SQL query. For example, given the question "List the top five students' performance", the piece of evidence \(k_1\) could state: "A student's performance is their average grade", thereby resolving the meaning of "performance" to the SQL construct `AVG(student\_courses.grade)`.

Using these definitions, we introduce the NL2SQL mapping function and our core formalizations.

\textbf{Definition 1. Mapping Function} \\
Given a database \textit{D} and external knowledge \textit{K}, we define a deterministic function \(f\) \footnote{As in [1], we leave the probabilistic definition to future work.} that maps a question \(q\) to a subset of \textit{S}, or to the empty set:
\[
f(q, D, K) \to S_q \subseteq S \cup \{\varnothing\}.
\]
The function \(f(q, D, K)\) returns the set of queries in \textit{S} that correctly answer \(q\) on \textit{D}, given the disambiguating information in \textit{K}. If no such query exists, \(f(q, D, K) = \varnothing\).

\textbf{Definition 2. Ambiguity} \\
A question \(q\) is \textit{ambiguous} with respect to \textit{D} and \textit{K} if \(f(q, D, K)\) yields two or more non-equivalent SQL queries. Formally, \(q\) is ambiguous if:
\[
\lvert f(q, D, K) \rvert \ge 2.
\]

\textbf{Definition 3. Unanswerability} \\
A question \(q\) is \textit{unanswerable} with respect to \textit{D} and \textit{K} if there exists no query in \textit{S} that can fulfill the request. Formally, \(q\) is unanswerable if:
\[
f(q, D, K) = \varnothing.
\]

\subsection{Taxonomy}
In this section, we explore our proposal for a taxonomy that focuses on the root cause of the \textit{ambiguity} and \textit{unanswerability} defined above. For each category, we bring some examples based on the 'University' database schema presented in Appendix \ref{schema}.

\subsubsection{Structural Ambiguity}
Structural Ambiguity occurs when the syntactic structure of a question permits multiple valid interpretations of the relationships among its components, resulting in distinct formal query mappings. Both scope and attachment ambiguities from \cite{Ambrosia} fall under this category.

\textit{Scope Ambiguity:} Arises from unclear quantifiers (e.g., each, every, all). For example, "What courses does each department offer?" can be interpreted collectively, referring to all departments together, or distributively, treating each department independently. These interpretations yield structurally different SQL queries: one aggregative, the other iterative.

\textit{Attachment Ambiguity:} Occurs from uncertainty in how a modifier attaches within the sentence. In "List the professors and the students in engineering," the phrase "in engineering" may modify only students or the entire conjunction professors and students, leading to distinct filtering conditions in the resulting query.

\subsubsection{Semantic Mapping Ambiguity}
Semantic Mapping Ambiguity occurs when a term or reference in the question can correspond to multiple schema elements, resulting in more than one plausible semantic grounding.

\textit{Lexical Overlap:} Arises when a natural language term maps to two or more schema attributes whose names share lexical stems, semantic roots, or refer to variants of the same concept. For example, in "List the emails of the students," the term "emails" could refer to \textit{students.personal\_email} or \textit{students.institutional\_email}—both columns share the concept of "email" but represent different types or variants. Similarly, "What is the date for all projects?" could refer to \textit{projects.start\_date}, \textit{projects.end\_date}, or \textit{projects.deadline}, where all column names contain "date" but refer to distinct temporal concepts. Each interpretation yields a distinct SQL mapping.

\textit{Entity Ambiguity:} Occurs when a natural language term or expression corresponds to attributes from multiple different entities or tables in the schema. The same concept exists across different tables representing different relational contexts. For instance, in "List the enrollment date of the students of the 'database' course," the expression "enrollment date" could refer either to the students' university enrollment (\textit{students.enrollment\_date}) or to the enrollment in the specific course (\textit{student\_courses.enrollment\_date}). The different entity interpretations require accessing different tables or following different join paths, as they refer to the same concept but in distinct relational contexts.

\subsubsection{Lexical Vagueness}
Lexical Vagueness arises when a question contains terms whose meaning lacks a precise or objective boundary, leading to indeterminate selection criteria during query generation. Such vagueness introduces variability that cannot be resolved solely from schema information, as it depends on the user's subjective understanding or context-specific conventions.

Common types of vague terms include:
\begin{itemize}
    \item \textbf{Temporal vagueness}: "recent," "old," "soon," "earlier" (e.g., "List recent courses")
    \item \textbf{Quantitative vagueness}: "many," "few," "high," "low," "significant" (e.g., "Show employees with high salaries")
    \item \textbf{Evaluative vagueness}: "popular," "good," "expensive," "important" (e.g., "Find popular products")
\end{itemize}

For instance, "List recent courses" is ambiguous with respect to the temporal threshold defining "recent"—it may refer to the last semester, the last academic year, or another undefined interval.

\subsubsection{Missing Schema Elements}
A question is unanswerable due to Missing Schema Elements when the database schema lacks the tables, columns, or relationships required to express the information requested.

\textit{Missing Entities or Attributes:} Occurs when key information is completely absent from the database schema—no table or column could represent the requested data. For example, "List the administrative staff in the engineering department" cannot be translated into SQL if the schema contains no table representing staff (assuming only students, professors, and courses are modeled). Similarly, "Show employee salaries for the IT department" would be unanswerable if an employees table exists but lacks a salary column. The missing elements are fundamental entities or attributes that would need to be added to the schema structure itself.

\textit{Missing Relationship:} Occurs when the necessary entities exist as separate tables but lack foreign keys, junction tables, or other relationship structures to connect them. For instance, "List the professors for the course 'database'" cannot be answered even though the schema includes both a \textit{professors} table and a \textit{courses} table, since no relationship (e.g., a \textit{teaches} junction table with \textit{professor\_id} and \textit{course\_id} foreign keys) links them to represent course assignments. Without these relationships, it is impossible to construct a SQL query that joins the relevant information together.

\subsubsection{Missing External Knowledge}
A question is unanswerable due to Missing External Knowledge when its interpretation depends on objective, domain-specific facts or policies not present in the database or knowledge base \textit{K}.

For instance, "List all the grades of the students for the course 'database' using alphabetic notation (A-F)" requires domain knowledge specifying the grade mapping (e.g., 'A+' → 31, 'A' → 30, etc.). Without this information, no valid SQL query can be formulated. Other examples include conversion formulas (e.g., "Convert all product prices to euros" requires exchange rates), classification rules (e.g., "Classify employees as junior, mid-level, or senior" requires experience thresholds), or institutional policies (e.g., "Find courses that satisfy the computer science degree requirements" requires degree requirement specifications).

\subsubsection{Missing User Knowledge}
A question is unanswerable due to Missing User Knowledge when its interpretation depends on objective but user-specific facts absent from the database or knowledge base \textit{K}.

For example, "List the students in my department" cannot be answered without knowledge of the user's department (e.g., User's department = 'engineering'). Similarly, "Show all projects I'm working on" requires knowing the user's identity or employee ID. These questions contain references to user context (e.g., "my," "our," "I") that require knowing who is asking the question and their specific attributes or affiliations. Without this user-specific information, the question cannot be properly resolved into a concrete SQL query.

\subsubsection{Conflicting Knowledge}
A question is ambiguous due to Conflicting Knowledge when the knowledge base \textit{K} contains multiple, mutually exclusive pieces of evidence for interpreting the same concept in the question. Crucially, the ambiguity arises not from the question wording itself being vague, but from having multiple documented, conflicting interpretations in the knowledge base. Each piece of evidence provides a valid but non-equivalent definition, leading to structurally different SQL queries. The user must specify which policy or evidence interpretation to follow.

For instance, consider the case where two evidences exist in \textit{K}: 
\begin{enumerate}
    \item "A student's performance is their average grade"
    \item "A student's performance is the average grade weighted by the course credits"
\end{enumerate}

Given the request "List the top five students' performance," the ambiguity stems from inconsistency within \textit{K}, not from vagueness in the term "performance" itself. The first evidence leads to a simple \texttt{AVG(grade)} calculation, while the second requires a weighted average computation. Both are valid interpretations according to different policies in \textit{K}.

\subsubsection{Improper Question}
A question is Improper when it is unrelated to the domain of either the database \textit{D} or the knowledge base \textit{K}.

This includes casual conversation or greetings (e.g., "Hello!"), questions requiring external reasoning or general knowledge not grounded in the database domain (e.g., "What is the meaning of life?"), or requests to perform actions rather than retrieve information (e.g., "Update my mailing address"). Such questions are fundamentally not database query requests and cannot be addressed through SQL generation.
